% =====================================================================================
% Camera-ready snippets for paper 462.  Paste into your existing source.
% Every  \todo{...}  is a number or sentence that must come from a REAL run -- search for
% "todo" before submitting and make sure none are left.  Do not fill them with estimates.
% Add to the preamble:   \usepackage{xcolor}  \usepackage{booktabs}
% =====================================================================================
\newcommand{\todo}[1]{{\color{red}\textbf{[#1]}}}

% ---- 0. Fix the broken reference "Fig. ??" in Sec. VI-A --------------------------------
% Put \label{fig:cm} INSIDE the confusion-matrix figure (after \caption) and write:
%   "... a macro-averaged F1 of 0.6789 (Fig.~\ref{fig:cm})."

% ---- 1. Sec. V-B  Dataset Preparation: add after the alias-mapping sentence ------------
\textbf{Duplicate prevention.} Before the split, every image was fingerprinted twice: a SHA-256
hash of the file bytes (exact duplicates) and a 64-bit difference hash with a Hamming-distance
threshold of 4 (re-encoded, resized or lightly cropped copies). The audit was run (i) over the raw
sources, to find images shared between the two datasets, and (ii) over the final train/validation/test
folders, to confirm that no image appears in more than one split. It found \todo{N} exact and
\todo{M} near-duplicate groups (\todo{P}\% of images), \todo{K} of which spanned both sources and
\todo{L} of which carried conflicting class labels. All but one copy of each group was removed
(\todo{R} images) before the stratified split, and the second audit reported \todo{0} cross-split
duplicates. The final dataset therefore contains \todo{8,770 or the new count} images.
% If the audit finds nothing, say so plainly: "The audit found no duplicate groups."

% ---- 2. Sec. VI: new subsection "Comparison with Other Architectures" -------------------
\subsection{Comparison with Other Architectures}
To test whether EfficientNetV2-S is a sound choice, four ImageNet-pretrained CNNs were trained on the
identical split with the identical recipe (\todo{30} epochs, batch size 32, AdamW with cosine
annealing, Mixup, label smoothing, seed 42) and evaluated on the same held-out test set
(Table~\ref{tab:model_comparison}). \todo{One or two sentences on what the table actually shows, e.g.
which model is most accurate, and the accuracy/latency trade-off. Only claim EfficientNetV2-S is best
if the numbers say so.}
% Paste the generated table: models/reports/model_comparison_table.tex  (or  \input{...})

% ---- 3. Sec. VI: new subsection "Portion-Size and Nutrient Estimation Error" -------------
\subsection{Portion-Size and Nutrient Estimation Error}
Nutrient values are looked up from a per-100\,g composition table rather than predicted, so two
error sources remain: the portion weight and the table values. \textbf{Portion.} \todo{n} servings of
\todo{f} dishes were weighed on a kitchen scale and photographed; the weight shown by the application
differed from the measured weight by an MAE of \todo{x}\,g (RMSE \todo{y}\,g, MAPE \todo{z}\%, mean bias
\todo{b}\,g). \textbf{Nutrients.} The stored per-100\,g values for \todo{F} foods were compared with
\todo{IFCT 2017 / USDA FoodData Central}; the MAE was \todo{a}\,kcal (MAPE \todo{c}\%) for energy,
\todo{..} for protein, \todo{..} for carbohydrate, \todo{..} for fat, \todo{..} for sugar and \todo{..}
for sodium. \textbf{End to end.} Combining both, the calorie estimate of a meal deviated from the
reference by \todo{MAE} kcal (\todo{MAPE}\%), of which \todo{table-only} kcal is attributable to the
table and \todo{portion-only} kcal to the portion estimate (Table~\ref{tab:errors}).
% Report n honestly; a small n is fine if you say it is small.

% ---- 4. Sec. VI: new subsection "Expert Validation of the Rules" ---------------------------
\subsection{Expert Validation of the Suitability Rules}
The thresholds and the verdicts they produce were reviewed by \todo{name/role, e.g. a registered
dietitian with N years of experience}. The reviewer checked every threshold (\todo{k} items) and gave an
independent verdict (ok / caution / avoid) for \todo{116} foods under each of the six conditions,
without seeing the system's verdict first; \todo{this blinding step only if it was really done}.
Agreement with the system was \todo{X}\% overall (Cohen's $\kappa=$ \todo{k}), ranging from
\todo{a}\% to \todo{b}\% across conditions, and \todo{Y}\% of the \todo{S} suggested substitutions were
judged acceptable. The reviewer recommended changing \todo{list of threshold changes}, which were
applied before the final results were generated. \todo{State the number of cases where the system said
"ok" and the expert said "avoid" -- this is the clinically important disagreement.}

% ---- 5. Replace the text under Fig. 4 (numbers recomputed after fixing a unit bug) --------
% The nutrition table stores values per 100 g, but earlier code scaled them as if they were per
% serving. With the fix, the worked example for Gulab Jamun (80 g, diabetes) is:
%   256 kcal, 3.2 g protein, 40.0 g carbohydrate, 9.6 g fat, 48 mg sodium, 32.0 g sugar.
% Verdict: avoid -- sugar 32 g vs 15 g limit, high glycaemic index, high-sugar food.
% (Carbohydrate, 40 g, is now BELOW the 45 g limit, so it is no longer a reason.)
% Alternatives are unchanged: Fresh Fruit Bowl, Plain Curd.  Re-capture the Fig. 4 screenshots from
% the fixed app and update: "Meal totals 256 kcal | 3.2 g protein | 40.0 g carbs | 9.6 g fat | 48 mg sodium | 32.0 g sugar".
% Also change "320 kcal" in the Fig. 4 caption/text to "256 kcal".

% ---- 6. Sec. VII Limitations: replace the portion sentence --------------------------------
% "Portion estimation uses the default serving unless a bounding box is available, so it does not
%  adapt to the amount actually served; its measured error is reported in Sec. VI-?."

% ---- 7. Sec. V-C: wording that matches the code ------------------------------------------
% "During the first epoch only the classification head was trained ..." is accurate for the updated
% training script (exact head parameters are selected). The earlier script also left EfficientNetV2's
% squeeze-excitation layers trainable in that epoch, which is why the final numbers must come from a
% run of the updated script.

% ---- 8. Table for Sec. VI "Portion-Size and Nutrient Estimation Error" (label tab:errors) ---
% Fill from models/reports/estimation_errors.json
\begin{table}[htbp]
\caption{Portion-size and nutrient-estimation error}
\label{tab:errors}
\centering\small
\begin{tabular}{lcc}
\toprule
\textbf{Quantity} & \textbf{MAE} & \textbf{MAPE (\%)} \\
\midrule
Portion weight (g), $n=$ \todo{n} & \todo{x} & \todo{z} \\
Energy per 100\,g (kcal), \todo{F} foods & \todo{a} & \todo{c} \\
Protein per 100\,g (g) & \todo{..} & \todo{..} \\
Carbohydrate per 100\,g (g) & \todo{..} & \todo{..} \\
Fat per 100\,g (g) & \todo{..} & \todo{..} \\
Sugar per 100\,g (g) & \todo{..} & \todo{..} \\
Sodium per 100\,g (mg) & \todo{..} & \todo{..} \\
Meal energy, end to end (kcal) & \todo{..} & \todo{..} \\
\bottomrule
\end{tabular}
\end{table}
